\section{模型评价与结论}

\subsection{四问结果与模型特点}

四个问题在相同物理状态方程与结算定义下递进。问题一的确定性LP解释储能如何在容量和效率约束下转移能量；问题二由五倍补购机制推导采购倾向，再用残差场景处理不确定性；问题三在合法交易时点按固定执行规则回放并筛选购电候选；问题四将同日价格残差与净负荷残差联合使用。最终配置为日前两组保留母模型、滚动两组在零时与日内均启用候选筛选，汇总如下。

\begin{table}[htbp]\centering\small
\caption{最终分问配置的334天开发期闭环评价结果}
\label{tab:annual_summary}
\begin{tabularx}{\linewidth}{LRRRR}\toprule
指标 & 问题2 & 问题3 & 问题4日前版 & 问题4滚动版\\\midrule
计划费用/万元 & 1350.80 & 1350.31 & 1428.01 & 1426.52\\
调整费用/万元 & 0.00 & 24.87 & 0.00 & 28.31\\
紧急费用/万元 & 50.48 & 11.97 & 51.40 & 12.19\\
实际总费用/万元 & 1401.28 & 1387.15 & 1479.41 & 1467.02\\
紧急电量/kWh & 85017.69 & 21029.90 & 80544.00 & 21293.41\\
紧急购电天数 & 113 & 119 & 106 & 116\\
紧急购电连续时段数 & 205 & 204 & 195 & 193\\
日总费用CVaR95/元 & 70731.74 & 65552.76 & 83044.31 & 77575.87\\
最高日费用/元 & 89306.82 & 67941.61 & 108080.37 & 81797.18\\
2月1日期初库存/kWh & 8875.19 & 8814.60 & 9251.54 & 9030.40\\
12月31日期末库存/kWh & 6240.32 & 6240.32 & 6269.86 & 6274.99\\
\bottomrule\end{tabularx}
\par\smallskip\parbox{\linewidth}{\footnotesize 问题2与问题4日前版沿用母模型；问题3与问题4滚动版在零时及日内均启用候选筛选。预测器与实时执行器沿用原冻结配置，筛选适用问题与时点依据开发期比较确定。实际费用不含候选评分的终端偏离项；CVaR95为实际日总账单的经验尾部均值。}
\end{table}


问题一费用为35126.95元，相对无储能节费26.90\%；问题二、三评价期主费用分别为1401.28万元、1387.15万元，动态电价日前版与滚动版分别为1479.41万元、1467.02万元。第三问与第四问滚动版的未筛选母模型费用分别为1397.58万元、1476.92万元，全时点筛选分别节省104339.41元和98994.66元，降幅为0.747\%和0.670\%。在同一问题四种配置末库存一致的条件下，全时点筛选取得本次比较中的最低账单。

这一收益伴随紧急补购增加：两滚动组分别由12647.76增至21029.90 kWh、由11950.68增至21293.41 kWh。计划费用的节省超过新增调整费用及紧急费用，并不表示所有运行指标同时改善。日前两组启用筛选反而分别增加0.161\%和0.428\%费用，因此保留母模型。信息来源、执行器和退款解释的原有对照均维持未筛选母模型配置，用于解释相应机制，不冒充最终增强算法的同条件消融。

主结果逐段满足供需平衡、储电上下界、功率限制、充放电互斥及日际连续性。未来实际值与未来发布预报的扰动没有改变已执行前缀。所有费用与时段分组由同一结果源生成，保留完整源代码、原模板及修订前版本，便于追溯图表中的每个数值。

\subsection{适用边界与可继续改进之处}

第一，只有一个年度且一月校准窗口较短；预测器与实时执行器沿用原冻结配置，筛选的适用问题和启用时点依据开发期结果确定，不能声称跨年份泛化或从未接触的独立测试。第二，预测与控制采用有限候选的坐标搜索，28场景、额外风险与备用关闭等设置保持不变；全时点筛选只在所比较的四种配置中费用最低，并未穷举全部可行策略。第三，候选回放沿用实际执行规则，部分改善了共享参考与实际控制的评估衔接，但回放固定当前合同至午夜，没有模拟后续尚未发布的预报及合同修订，不是完整多阶段随机最优控制。第四，终端偏离罚项仅用于优化和候选评分，不计入实际账单；模型未加入题目以外的网络限额、设备寿命退化和交易约束。

扩大候选集只能保证相同情景、初态和评分规则下的最小估计费用不增加，不能保证未来实际费用下降。本次保留具有消融依据的滚动采购增强，不再据开发期结果叠加风险、备用或新的实时控制器；其适用性仍需在其他未参与选型的时段或年度检验。
